\scanpage{222}
\deleted{\illegible{1}}
\diagram{\begin{tikzcd}[ampersand replacement=\&,column sep=small]
X^{(0)}_{\deleted{\unsure{a}}}\arrow[d] \& \deleted{X}\arrow[d] \\
S^{(0)}_{\deleted{\unsure{0|e}}}=\operatorname{Spec}(\mathbb Z[t])\deleted{\unsure{k}} \& S\arrow[l]\\
\& \tau\in\Gamma(S,\underline{\mathcal O}_S)
\end{tikzcd}}
\[
\left(\begin{aligned}
&A[x,y]/xy(x+y+t)\\
&A[x,y,z]/xy(x+y+\tau z)\deleted{\unsure{=A}}=\mathscr S
\end{aligned}\right.
\]
\[
\deleted{xy(x+y+tz)}
\]
\diagram{$\begin{array}{ccc}
(x)&(y)&(x+y+tz)\\
\unsure{\mathfrak p|p}_1&\mathfrak p_2&\mathfrak p_{3}\\
\underline{\unsure{\mathcal J|\mathcal I|I}}_1&\underline{\mathcal J}_2&\underline{\mathcal J}_3\\
X_1&X_2&X_3
\end{array}$}

\uline{Lemme 1} $X$ est plat sur $S$

Il suffit de prouver que $X^{(0)}$ est plat sur $S^{(0)}$, donc à fortiori \deleted{\unsure{de prouver que}} $\unsure{k|\mathbb Z}[x,y,z,t]/xy(x+y+tz)$. est plat sur $k[t]$ ; \deleted{donc} i.e. sans torsion. Or supposons que $\varphi\in k[t]$, \inserted{$\varphi\ne0$}, $F\in k[x,y,z,t]$ tels que $\varphi F\in xy(x+y+tz)$, donc
\[
\begin{aligned}
\varphi F&\in(x)\\
\varphi F&\in(y)\\
\varphi F&\in(x+y+tz)
\end{aligned}
\]
d'où Comme \deleted{$k[x,y,z,t]$} les idéaux \deleted{\unsure{$(x),(y),(x+y+tz)$}} $\mathfrak p_1^{(0)}\unsure{:|=}(x)$, $\mathfrak p_2^{(0)}\unsure{:|=}(y)$, $\mathfrak p_3^{0}=(x+y+tz)$ sont \margin{\uline{Lemme \unsure{1|2}}} \uline{premiers}, et $k[t]\longrightarrow k[x,y,z,t]/\mathfrak p_i^{(0)}$ est \uline{injectif} \margin{\deleted{\uline{Lemme}}} la relation $F\in\deleted{\unsure{(x), F\in(y)}}\mathfrak p_i^{(0)}$ pour \unsure{tt} $i$ \deleted{\unsure{pour|$=1,2,3$}} d'où $F\in\mathfrak p=xy(x+y+tz)$ à cause de

\uline{Lemme \unsure{2|3|0}} On a $\mathfrak p_1\cap\mathfrak p_2\cap\mathfrak p_3=0$ i.e. $(x)\cap(y)\cap(x+y+zt)=(xy\unsure{(x+y+|x+y+}$
% Right edge cuts through the last factor.
% archive sheet number: 221
